% What is the low-field mobility
Experiments conducted thereafter have showed mobility extending between $\sim 3000 - 12000$ cm$^2$ V$^{-1}$s$^{-1}$ for carrier concentrations $n\simeq1 - 0.1 \times 10^{12}$ cm$^{-2}$, with
low-field field-effect mobility\footnote[2]{This differs from the Drude ($\mu_{\textrm{Drude}}=\sigma/ne$) and Hall ($\mu_\textrm{Hall}=-\sigma_e R_H$) mobilities.} ($\mu = (1/C_{\textrm{g}})(\partial \sigma/\partial V_{\textrm{gs}})$)
that is constant over a quite wide range of $n$ (yielding in this case $\mu_h\simeq1900$ cm$^2$ V$^{-1}$ s$^{-1}$ and $\mu_e\simeq1800$ cm$^2$ V$^{-1}$ s$^{-1}$ for the low $n$ region)~\cite{Meric2008}.


Later in 2008~\cite{bolotin2008ultrahigh} and 2013~\cite{wang2013one}, higher mobility values were reported:  at 240 K and $\simeq1,000,000$ cm$^{2}$V$^{-1}$s$^{-1}$ at 4 K. These higher values were attributed to the use of suspended graphene as opposed to graphene on a substrate, where surface effects inevitably affect carrier dynamics. Moreover, mobilities in the order of $\simeq120,000$ cm$^{2}$V$^{-1}$s$^{-1}$ obtained at room temperature with graphene deposited on a single layer of hexagonal boron nitride~\cite{dean2010boron} and sandwiched between two h-BN layers~\cite{wang2013one} indicate that this flat and inert material can serve as an excellent passivation film for graphene devices.


It is worthwhile to consider the physics of the metal-semiconductor junction (M/S) in order to understand the physical origin of this figure of merit. The thermoionic emission theory and the exclusion of other types of conduction mechanisms (thermoionic field emission, generation/recombination, electron diffusion, injection of holes, and leakage current) may reasonably reproduce the $I$--$V_\textrm{SD}$ characteristics of such junction. Within this model, the off current is proportional to the barrier that the electrons must overcome in order to move from the semiconductor to the metal:

\begin{equation}
    I_\textrm{off}\propto\exp{(e\beta\phi_B)}.
\end{equation}


Back-gate devices suffered from parasitic capacitances, making further integration with other components impossible despite their usefulness for proof-of-concept purposes. Alternative geometries have since been investigated to increase carrier mobility, decrease contact resistances and parasitic capacitances, and eliminate the Schottky barrier between graphene and semiconductor.

graphene-based  have not been adopted by industry and have remained a purely academic pursuit.


The Manchester group has recently presented a novel approach to graphene FETs with vertical tunnel current flow between parallel graphene layers acting as source and drain (Britnell et al., 2012). Several devices are constructed with single-crystal layers of graphene and hexagonal boron nitride, and additional devices are constructed with tunnel barrier MoS2. The source and drain are separated by a tunnel barrier, with the sandwich placed horizontally on the gate electrode's oxidised Si wafer. This FET device lacks a conventional "channel," as the forward device current flows vertically from the source to the drain via tunnelling. The gate voltage modifies both the carrier concentrations and the tunnelling probability across the single crystal tunnel barrier. Remarkably, micromechanical cleavage has made it possible to create nearly ideal tunnel barriers and I(V) curves for tunnelling between two graphene monolayer electrodes. The tunnel barrier consists of a few layers of mechanically exfoliated BN or MoS2 film.

For , however, full-speed operation is attained by operating in the (drain) current saturation regime.

In \gls{RF} applications, however, the operational frequency at which the carriers run generates heat. Thus, two additional characteristics are necessary: first, a high thermal conductivity, and second, a low heat transfer with the substrate. Here, trade-offs are required.

In 2007, Lemme and coworkers used a stack of \ce{SiO2} (20 nm), Ti (10 nm), and Au (100 nm) as the top gate for modulating the current inside a large-area graphene channel transistor obtained by the exfoliation method, heralding the introduction of vertical graphene transistors~\cite{lemme2007graphene},.

Whilst the gate characteristics exhibit a similar trend of the characteristics (the top gate reduces the back-gate current modulation, lowering the drain current by one order of magnitude) and the hysteresis, weak current saturation at high drain-source bias have been reported~\cite{moon2009epitaxial}.
This effect seem due to a partial charge pinning at the contacts at high back-gate voltages~\cite{di2011charge} and carrier drift velocity eventually saturating due to optical-phonon scattering~\cite{meric2008current}. Top-gated graphene \glspl{MOSFET} with CVD and epitaxial graphene have been also made~\cite{kedzierski2008epitaxial, lin2010100}; \ce{SiO2}, \ce{Al2O3}, and \ce{HfO2} have been  used as top-gate dielectric.

However, despite the improvements obtained with a top gate, most graphene \gls{MOSFET} with large-area graphene channels show on-off switching ratios of less than 10 at room temperature~\cite{schwierz2010graphene}, much too little for complex logic circuits.

In 2012, Mehr \textit{et al}.\ proposed a vertical arrangement of emitter, base, and collector similar to the original vacuum triode, where graphene plays the base role and electrons travel vertically rather than transversely.
When the device is in the ON state, the electrons tunnel with a Fowler-Nordheim tunnelling mechanism:

\begin{equation}
    I\propto\qty(\frac{V_\textrm{SD}}{d})^2~\exp{-\frac{\alpha~d}{V_{sd}}},
\end{equation}

through the emitter--base layer (of thickness $d$); ON/OFF ratios of the order of $10^4$ have been demonstrated using graphene between two h-BN layers~\cite{britnell2012field} and using \ce{WeS2}--graphene~\cite{georgiou2013vertical} heterostructures.




In order to be used in logic applications graphene should have a highly conductive on-state and an off-state in which no source-drain $I_\textrm{SD}$ current flows. This requirement may be partially filled by requiring to have a variable bandgap.For \gls{RF} applications, however, full-speed operation is attained by operating in the (drain) current saturation regime. Moreover, the trend of \gls{IC} scaling is frequently contrasted with chip-level concerns such as power densities (now on the order of 1000 W/cm$^2$), heat production, and chip temperatures approaching levels that impede stable operations.

The first prototypes of graphene-based \glspl{FET} were made through back-gate contact. In this geometry, tipically the gate is a Si/\ce{SiO2} substrate on top of which graphene (\gls{CVD}, single-layer, epitaxial) is placed thorugh a transfer method (Fig.~\ref{fig:intro:graphene_device}).

In CMOS logic, charge carriers are driven by the electric field  $\CMcal{E}_\textrm{sd}$ between source and drain, resulting in four potential sources of scattering: electron--electron, electron-phonon (lattice vibrations), material interfaces, and impurity atoms. The first two are more foundational, whereas the last two are connected to the type of interfaces formed between graphene and its support material. The predominant process through which charge carriers lose net energy is scattering with phonons, resulting in Joule heating of the lattice. In \gls{RF} applications, however, the operational frequency at which the carriers run generates heat. Thus, two additional characteristics are necessary: first, a high thermal conductivity, and second, a low heat transfer with the substrate. Here, trade-offs are required.

The 2020 International Roadmap for Devices and Systems, which followed the 2017 International Roadmap for Semiconductors, places a significant emphasis on graphene in its study \textit{beyond CMOS}, which investigates technologies and semiconductor materials that will supplement and/or replace Si. Since carrier mobility is a well acknowledged characteristic for addressing carrier transport in nanoscale \gls{FET}, the race to scale the channel continues to motivate significant efforts to introduce high-mobility channel materials.

Indeed, the majority of the initial interest in adopting graphene as a channel material in standard FET technology was sparked by the extraordinary mobility values reported since its discovery (Tab.~\ref{tlc}). \glspl{FET}~\cite{schwierz2010nanometer}. At \SI{290}{\kelvin}, graphene's intrinsic mobility $\mu_e^{\textrm{in}}$ is extraordinarily high. The experimental data reported in 2004 revealed mobilities close to $\mu_e^\textrm{in}\simeq10,000$ cm$^{2}$V$^{-1}$s$^{-1}$~\cite{novoselov2004electric}.
Experiments conducted thereafter have showed mobility extending between $\mu^{\text{in}}\sim 3000 - 12000$ cm$^2$ V$^{-1}$s$^{-1}$ for carrier concentrations $n\simeq1 - 0.1 \times 10^{12}$ cm$^{-2}$, with
low-field field-effect mobility\footnote[2]{This differs from the Drude ($\mu_{\textrm{Drude}}=\sigma/ne$) and Hall ($\mu_\textrm{Hall}=-\sigma_e R_H$) mobilities.} ($\mu = (1/C_{\textrm{g}})(\partial \sigma/\partial V_{\textrm{gs}})$)
that is constant over a quite wide range of $n$ (yielding in this case $\mu_h\simeq1900$ cm$^2$ V$^{-1}$ s$^{-1}$ and $\mu_e\simeq1800$ cm$^2$ V$^{-1}$ s$^{-1}$ for the low $n$ region)~\cite{Meric2008}.




Later in 2008~\cite{bolotin2008ultrahigh} and 2013~\cite{wang2013one}, higher mobility values were reported: $\simeq120,000$ cm$^{2}$V$^{-1}$s$^{-1}$ at 240 K and $\simeq1,000,000$ cm$^{2}$V$^{-1}$s$^{-1}$ at 4 K. These higher values were attributed to the use of suspended graphene as opposed to graphene on a substrate, where surface effects inevitably affect carrier dynamics. Moreover, mobilities in the order of $\simeq120,000$ cm$^{2}$V$^{-1}$s$^{-1}$ obtained at room temperature with graphene deposited on a single layer of hexagonal boron nitride~\cite{dean2010boron} and sandwiched between two h-BN layers~\cite{wang2013one} indicate that this flat and inert material can serve as an excellent passivation film for graphene devices.

\co{Ion/Ioff ratio}
Another figure of merit one has to take into account when designing \glspl{FET} is the $I_\text{on}/I_\text{off}$. On-off ratio is the ratio of the current in the on-state to the off-state current when no gate is applied. A high on-off ratio corresponds to a low leakage current, resulting in an enhanced device performance. Due to the short channel effects, the off-current of a \glspl{MOSFET} increases as the channel length scales down. While the scaling of a graphene-based transistor could be, in principle, be robust against this effect, Klein tunnelling often represents the major hurdle when graphene is employed~\cite{schwierz2013graphene}.

It is worthwhile to consider the physics of the metal-semiconductor junction (M/S) in order to understand the physical origin of this figure of merit. The thermoionic emission theory and the exclusion of other types of conduction mechanisms (thermoionic field emission, generation/recombination, electron diffusion, injection of holes, and leakage current) may reasonably reproduce the $I$--$V_\textrm{SD}$ characteristics of such junction. Within this model, the off current is proportional to the barrier that the electrons must overcome in order to move from the semiconductor to the metal:

\begin{equation}
    I_\textrm{off}\propto\exp{(e\beta\phi_B)}.
\end{equation}

where $e$ is the electron charge, $\phi_B$ is the Schottky barrier height, and $\beta=k_BT$. A pictorial representation is given in Fig.~\ref{fig:intro:graphene_device}
Assuming that the work functions of graphene and the metal are identical, the barrier height simply becomes $\phi_B = E_g/2$, where $E_g$ is the bandgap energy of the semiconductor, and the on-off ratio may be estimated from the simple relation ~\cite{kim2011role}:

\begin{equation}\label{eq:bandgap_current_ratio}
    \frac{I_\textrm{on}}{I_\textrm{off}}\propto\exp{\frac{E_g}{2k_BT}},
\end{equation}

where the thermal energy $k_BT$ and the bandgap are both expressed in meV. \Gls{CMOS} bandgap values should fall between 300 and 500, according to a number of scholarly references~\cite{schwierz2010graphene, reddy2011graphene}.

Back-gate devices suffered from parasitic capacitances, making further integration with other components impossible despite their usefulness for proof-of-concept purposes. Alternative geometries have since been investigated to increase carrier mobility, decrease contact resistances and parasitic capacitances, and eliminate the Schottky barrier between graphene and semiconductor.

In 2007, Lemme and coworkers used a stack of \ce{SiO2} (20 nm), Ti (10 nm), and Au (100 nm) as the top gate for modulating the current inside a large-area graphene channel transistor obtained by the exfoliation method, heralding the introduction of vertical graphene transistors~\cite{lemme2007graphene},.

Whilst the gate characteristics exhibit a similar trend of the characteristics (the top gate reduces the back-gate current modulation, lowering the drain current by one order of magnitude) and the hysteresis, weak current saturation at high drain-source bias have been reported~\cite{moon2009epitaxial}.
This effect seem due to a partial charge pinning at the contacts at high back-gate voltages~\cite{di2011charge} and carrier drift velocity eventually saturating due to optical-phonon scattering~\cite{meric2008current}. Top-gated graphene \glspl{MOSFET} with CVD and epitaxial graphene have been also made~\cite{kedzierski2008epitaxial, lin2010100}; \ce{SiO2}, \ce{Al2O3}, and \ce{HfO2} have been  used as top-gate dielectric.

However, despite the improvements obtained with a top gate, most graphene \gls{MOSFET} with large-area graphene channels show on-off switching ratios of less than 10 at room temperature~\cite{schwierz2010graphene}, much too little for complex logic circuits.

In 2012, Mehr \textit{et al}.\ proposed a vertical arrangement of emitter, base, and collector similar to the original vacuum triode, where graphene plays the base role and electrons travel vertically rather than transversely.
When the device is in the ON state, the electrons tunnel with a Fowler-Nordheim tunnelling mechanism:

\begin{equation}
    I\propto\qty(\frac{V_\textrm{SD}}{d})^2~\exp{-\frac{\alpha~d}{V_{sd}}},
\end{equation}

through the emitter--base layer (of thickness $d$); ON/OFF ratios of the order of $10^4$ have been demonstrated using graphene between two h-BN layers~\cite{britnell2012field} and using \ce{WeS2}--graphene~\cite{georgiou2013vertical} heterostructures.

Solution methods have also been explored, in which graphene is dropcasted on top of a semiconducting wafer, and then spin-coated~\cite{liu2010simple}.
The main difficulties arise when attempting to quantify the structure and the quality of the deposited graphene, and not many methods can guarantee a good dispersion in organic solvents~\cite{sun2011graphene}.
Simpler ways to obtain graphene oxide in solution have been reported~\cite{van2016effect}, but once deposited, graphene oxide needs to be reduced for obtaining pure graphene -- low mobility values (up to, at best, 1000 \si{\square\centi\metre\per\volt\per\second}) caused by defects during the fabrication steps have discouraged further investigations along this direction.

Yet in spite of the very different dc behavior compared to the most used semiconductors in digital electroncis, graphene FETs have reported good performance in terms of high-frequency applications. This comparison is made by measuring each system’s frequency response, its output signal strength when it is given a particular input signal: features such as ballistic transport, linear current-voltage characteristics, high transconductance, and very high ($\simeq$ \SI{1e8}{\ampere\per\square\centi\metre}) sustainable currents are all ideal features for analog/RF applications~\cite{meric2008rf, meric2008current}.



Late in 2008, it was claimed that the first device with a cutoff frequency in the gigahertz region was farbicated using exfoliated graphene placed on a high-resistivity silicon substrate, with a 30-nm-thick top-gate dielectric made of \ce{HfO2} and a 500-nm-long gate. Despite a reported cutoff frequency of about $f_\textrm{T}\simeq\SI{14.7}{\giga\hertz}$, it was stated that excessive contact resistances and gate length were the primary causes of low ($f_\textrm{max}<\SI{1}{\giga\hertz}$) maximum frequency performance. A 350-nm gate utilised Pd/Au contacts and 12-nm-thick aluminium oxide the next year. While the cutoff frequency improved up to 50 GHz, the maximum frequency only increased marginally to 2 GHz.

The first device with cutoff frequency $f_\textrm{T}$ in the GHz range was reported in late 2008, and was farbicated by exfoliated graphene deposited on top of high-resistivity \ce{Si} substrate, with 30-nm-thick top-gate \ce{HfO2} dielectric and 500 nm of gate length~\cite{meric2008rf}.
Despite a measured $f_\textrm{T}\simeq 14.7$ GHz, it was argued that large contact resistances and gate length were the main responsible of poor perfomance $f_\textrm{max} < 1$ GHz.

The following year, Pd/Au contacts and 12-nm-thick \ce{Al2O3} were employed in a 350-nm-gate~\cite{lin2009development}. While the cutoff frequency improved up to $f_\textrm{T}\simeq 50$ GHz, little gains in terms of $f_\textrm{max}$ were obtained of only 2 GHz. In general, limited performance in $f_\textrm{max}$
represents a major issue as for most RF applications $f_\textrm{T}$ is less important than $f_\textrm{max}$.

It was only after two years that an increase of one order of magnitude both in the cutoff and the maximum frequency of oscillation was demosntrated~\cite{lin2010100}. The device was realised with epitaxially-grown on silicon carbide (\ce{SiC}) graphene \gls{MOSFET} having 240 nm of gate length, with Ti/Pd/Au for contacts and 10-nm-thick \ce{HfO2} of insulating layer on top gates. Highest frequency as well as $f_\textrm{T}\simeq 100$ GHz and $f_\textrm{max}\simeq 14$ GHz were measured,
exceeding any device then present. The authors also argued that since $f_\textrm{max}$ is strongly correlated to device layout and parasitic resistances originating from the metallic contacts, viable ways to improve this parameter would be a partial passivation of the metal contacts and a multi-finger gate layout

These perfomance were soon followed later that year by a 144-nm-length gate device with $f_T \simeq 300$ GHz using mechanically exfoliated graphene and \ce{Co2Si}–\ce{Al2O3} nanowire as self-aligned multi top-gates~\cite{liao2010high}. This result was also corrobated by comparing it with the projected
value using the relation that links dc transconductance and gate capacitance with the cutoff frequency when total input parasitic capacitance is negligible, namely $f_\textrm{T} = g_\textrm{m}/(2\pi C_\textrm{g})$~\cite{sze2021physics}. The authors did not report the power $G^{1/2}_\textrm{max}$ but only the current $\abs{h_{21}}$ gain, thus hampering an evaluation of $f_\textrm{max}$,
but in 2012 a similar proposal was advanced by Cheng \textit{et al.}, where Au/Ti/\ce{Al2O3} gate stacks were employed in self-aligned transistors~\cite{cheng2012high}. Despite the improvements reported ($f_\textrm{T}\simeq 427$ GHz for a channel length of 67 nm) respect to CVD graphene transistors with comparable channel length, $f_\textrm{max}$ decresead down to 8 GHz when the length was shrunk down to 46 nm, showing the
key role played by gate and parasitic resistances at high frequency.

To further improve the power gain, works on the saturation behavior of short channel graphene transistor were also conducted by IBM in 2012, by using CVD graphene grown on a Cu foil and then transferred on diamond-like and SiC substrates~\cite{wu2012state}. These are appealing because high-quality
wafers are available, thus making a scale up of graphene arrays relatively straightforward. Since the work function of graphene is $\phi^\textrm{g} = 4.5$ eV and $\phi^\textrm{4H-\ce{SiC}} = 3.7$ eV, this means that $\Delta\phi^\textrm{g--4H-SiC} = 0.8$ eV, thus $n$-doping graphene.
Very cleverly, the authors then used plasma-enhanced CVD silicon nitride as an intrisic p-type gate dielectric, counterbalancing the n-doped graphene channels. Encouraing performance at high speed were not only measured in good cutoff frequencies ($f_\textrm{T}\simeq 350$ GHz), but also in a factor of 2 improvement
of $f_\textrm{max}$ ($\simeq 40$ GHz).
To the best of my knowledge, the current graphene transistor power gain record of about 70 GHz was obtained in 2013 by using (000-1) carbon-terminated surface of SiC instead of the (0001) Si one, indicating the contact resistance $R_\textrm{c}$ between Pd/Au contacts and graphene as key responsible for the high $f_\textrm{max}$.

While these results have aligned graphene MOSFET with the main solid-state competitors -- such as GaAS pHEMT, InP HEMT and Si MOSFET -- in terms of $f_\textrm{T}$, the perfomance is terms of $f_\textrm{max}$ are off by about one~\cite{tan1990InGaAs, post2006RF} or two orders of magnitude~\cite{lai2007sub}. In literature, one of the often cited fundamental reason
for this behaviour is the poor drain saturation current in graphene, which ultimately depends on the missing bandgap of the channel. Graphene-derived materials such as nanoribbons and carbon nanotubes MOSFET could then, given their intrisic semiconducting nature, represent a possible solution to this problem.



